\scanpage{321}

\mspage{60}
% archive sheet number: 320

est une flèche de $a$ vers $b$. Dans le cas général, les chemins de $a\unsure{)|,}$ vers $b\unsure{)}$ peuvent être vus comme une sorte de ``flèches généralisées'' allant de $a$ vers $b$, dont la composition \deleted{est la} serait la composition \unsure{évidente}, et les identités les ``chemins vides'' relatifs aux $a\in\operatorname{Ob}X$. \deleted{\unsure{D|Il}}\unsure{On construit|On obtient} d'ailleurs bien ainsi une catégorie, ayant $=$s objets que $X$, mais où les
\[
\operatorname{Hom}_{\mathrm{ch}}(a,b)=\mathrm{Ch}(a,b)=\operatorname{Ob}\underline{\mathrm{Ch}}(a,b)
\]
sont prohibitivement gros. \deleted{\unsure{Il}} Cela s'exprime \inserted{notamment} par le fait que les seules flèches inversibles dans $X_{\mathrm{ch}}$ sont les identités! On n'a pas de foncteur naturel entre $X_{\mathrm{ch}}$ et $X$, \deleted{\unsure{dans le}}\inserted{qui} soit l'identité sur les objets. Car à un chemin $a\overset{c}{\unsure{\text{ --- }|\longrightarrow}}b$ on ne peut associer de $X$-morphisme $a\text{ --- }b$ (sauf si $c$ de type $\Delta_+$), et à un $X$-morphisme $a\overset{u}{\text{ --- }}b$ on peut associer un chemin \inserted{$c(u)$} de type $\Delta_+$, mais \deleted{\unsure{ce}} \inserted{si $a\overset{u}{\text{ --- }}b\overset{v}{\text{ --- }}c$,} $c(vu)$ est \inserted{encore} de type $\Delta_+$ alors que $c(v)c(u)$ est de type $\Delta_+^2$, ils ne sont pas égaux. Par contre, il y a une flèche \inserted{naturelle} entre ces deux \deleted{\unsure{chemins}} éléments de $\underline{\mathrm{Ch}}(a,\unsure{c|b})$, flèche au dessus de \deleted{l'\unsure{hom. entre}} l'unique hom.
\[
\Delta_+\xrightarrow{f_{+,2}}\Delta_+^2
\]
